\documentclass[11pt]{article}
\usepackage[margin=1in]{geometry}
\usepackage{amsmath,amssymb,amsthm,mathtools}
\usepackage{hyperref}
\usepackage{enumitem}
\usepackage{booktabs}
\usepackage{array}

\newtheorem{theorem}{Theorem}[section]
\newtheorem{corollary}[theorem]{Corollary}
\newtheorem{lemma}[theorem]{Lemma}
\newtheorem{conjecture}[theorem]{Conjecture}
\newtheorem*{definition}{Definition}
\newtheorem*{remark}{Remark}

\newcommand{\Q}{\mathbb{Q}}
\newcommand{\Z}{\mathbb{Z}}
\newcommand{\N}{\mathbb{N}}
\newcommand{\repo}[1]{\text{\texttt{#1}}}
\newcommand{\cw}{c_W}

\hypersetup{colorlinks=true, linkcolor=blue, citecolor=blue, urlcolor=blue}

\title{Beyond Sturmian Words in the 3x+1 Periodicity Conjecture:\\
A Periodic-Approximation and Arithmetic-Height Bridge}
\author{Elias De Jesús\\
\textit{Independent Researcher}\\
ORCID: \href{https://orcid.org/0009-0007-0190-9143}{0009-0007-0190-9143}}
\date{September 2026}

\begin{document}
\sloppy
\maketitle

\begin{abstract}
Let $T:\Z_2\to\Z_2$ be the $3x+1$ map and $\Phi$ the Bernstein--Lagarias conjugacy map sending a parity vector $v\in\{0,1\}^\N$ to the 2-adic integer it codes. Lagarias's Periodicity Conjecture asserts that $\Phi(v)\in\Q$ only if $v$ is eventually periodic; it is open in general. A companion unpublished paper by the present author proves it for every Sturmian parity word, by two independent mechanisms (a transported counting bound, and an effective Liouville-style argument using an initial square). This note consolidates a subsequent research program that pushes the periodic-approximation/arithmetic-height mechanism of the second route past the Sturmian case, into symbolic classes the counting route cannot reach. The main proved result: for \emph{every} irrational slope with bounded partial quotients, every complementary symmetric Rote sequence (factor complexity exactly $2n$) built on that slope's intercept-$0$ mechanical word --- either seed, any shift --- has irrational $\Phi$-value. We give the full proof architecture, including a structural fact (an odd-weight transfer route's doubled root always has density exactly $1/2$) that collapses an earlier, over-conservative repetition-exponent requirement down to the true one. We document, rather than hide, four working hypotheses this program formed and then disproved by direct adversarial testing. We map --- but do not resolve --- the obstruction that appears once the intercept-$0$ restriction is lifted, reducing the general problem to one precisely stated, computationally well-supported but unproved margin lemma. A late-stage exploratory probe into a conjectural symbolic-symmetry/arithmetic-admissibility tradeoff is reported honestly as inconclusive, not as a result. \textbf{This note does not prove, and does not claim progress toward proving, the Periodicity Conjecture in general}; a structural fact recorded in the underlying research record shows no purely repetition-based method, however extended, can reach every aperiodic word.
\end{abstract}

\tableofcontents

\section{The Periodicity Conjecture and the Conjugacy Map}
\label{sec:conjecture}

Let $T:\Z_2\to\Z_2$ be the $3x+1$ map in its standard form,
\[
T(x)=\begin{cases} x/2, & x\equiv 0 \pmod 2,\\ (3x+1)/2, & x\equiv 1\pmod 2.\end{cases}
\]
The \emph{parity vector} of $x\in\Z_2$ is $v(x)=(v_0,v_1,\dots)\in\{0,1\}^\N$ with $v_i\equiv T^i(x)\pmod 2$. Bernstein and Lagarias showed that $x\mapsto v(x)$ is a bijection of $\Z_2$ onto $\{0,1\}^\N$, and in fact a \emph{2-adic isometry} \cite{BernsteinLagarias1996,Bernstein1994}: identifying $\{0,1\}^\N$ with $\Z_2$ through binary expansions,
\[
v_2\bigl(\Phi(a)-\Phi(b)\bigr) = \operatorname{lcp}(a,b),
\]
where $\Phi$ is the inverse map (so $\Phi(v)$ is the unique 2-adic integer with parity vector $v$), and $\operatorname{lcp}$ is longest-common-prefix length. Conjugating $T$ by $\Phi$ turns it into the shift map on $\{0,1\}^\N$. This isometry property is the single fact the entire bridge in this note exploits: two sequences agreeing on a long prefix must have 2-adically close $\Phi$-images.

Rationality of $\Phi(v)$ corresponds to eventual periodicity of $v$ in one direction only: if $v$ is eventually periodic then $\Phi(v)\in\Q$. The converse is:

\begin{quote}
\textbf{Lagarias's Periodicity Conjecture} \cite[\S2]{Lagarias1985}. No aperiodic parity vector has a rational $\Phi$-value.
\end{quote}
(See also López and Stoll's continuous-setting formulation of the same conjecture, \cite{LopezStoll2021}.)

This conjecture is open, and \emph{very little is known unconditionally in the converse direction} --- a fact worth stating plainly, since it frames everything that follows. The conjecture quantifies over \emph{all} $v\in\{0,1\}^\N$ that fail to be eventually periodic: an uncountable family with no assumed combinatorial structure.

\paragraph{Full conjecture versus the classes treated here.} This note --- and the research program it consolidates --- does not address the full conjecture. It proves $\Phi(v)\notin\Q$ for two structurally defined sub-families of aperiodic words (Sturmian words, proved in the companion foundation paper; and complementary symmetric Rote sequences over bounded-type slopes at intercept $0$, proved here), and precisely maps, without resolving, one further natural extension (Rote sequences at arbitrary intercepts). \S\ref{sec:open} records a structural fact showing that no extension of this program's method, however far pushed, can by itself close the full conjecture. Nothing in this note should be read as narrowing that gap beyond what is explicitly proved.

\section{Two Routes to Irrationality}
\label{sec:tworoutes}

The foundation paper and this program's extension of it use two logically independent mechanisms to force $\Phi(v)\notin\Q$.

\subsection{Route A --- Counting}
\label{sec:routeA}

Let $p_v(L)$ denote the number of distinct length-$L$ factors of $v$ (its \emph{factor complexity}). Dubickas proved in 2009 \cite{Dubickas2009} that a positive-integer $3x+1$ trajectory tending to infinity has parity-word complexity $p(L) > 1.70951129\,L$ eventually; the foundation paper transports this to every rational 2-adic integer, since $T$ preserves a fixed odd denominator. The exact threshold constant is
\[
\frac{1}{\log_2(3/2)} = 1.709511\ldots,
\]
and the transported statement is: if $v$ is not eventually periodic and $\Phi(v)\in\Q$, then $\liminf_L p_v(L)/L \ge 1/\log_2(3/2)$. Contrapositively: \emph{any} aperiodic $v$ with $\liminf_L p_v(L)/L < 1.70951\ldots$ has $\Phi(v)\notin\Q$, unconditionally, with no repetition or discrepancy hypothesis whatsoever.

This is a purely counting-based exclusion, and it is exactly as strong as its complexity bound: it says nothing once a word's complexity crosses the threshold. Sturmian words ($p(n)=n+1$) sit far below it. Complementary symmetric Rote sequences, the next class this program treats, have $p(n)=2n$ --- strictly above $1.70951\ldots\,n$ for every $n\ge 2$ --- and are therefore \emph{entirely outside Route A's reach}. This is precisely where Route A stops, and why a second, structurally different mechanism is needed to make any further progress.

\subsection{Route B --- Periodic Approximation and Arithmetic Height}
\label{sec:routeB}

This is the central mechanism of this note. Informally, the chain of implications is: symbolic periodic approximation $\to$ exact 2-adic agreement $\to$ divisibility $\to$ rational periodic approximant $\to$ Archimedean height $\to$ contradiction.

Concretely: if $v$ agrees with a periodic word $W^\infty$ on a long prefix, the 2-adic isometry forces $\Phi(v)$ and the \emph{rational number} $\Phi(W^\infty)$ to be 2-adically close to a precisely quantified depth. If $\Phi(v)$ were itself rational, this closeness becomes an exact integer divisibility statement between two specific integers built from $\Phi(v)$'s numerator/denominator and $W$'s combinatorial data. When the periodic approximation is long and clean enough relative to the arithmetic size (\emph{height}) of the numbers involved, that divisibility is impossible unless the two rationals coincide --- and they were assumed not to. The whole of \S\ref{sec:abstract}--\S\ref{sec:repheight} makes this precise.

\section{The Abstract Bridge Theorem}
\label{sec:abstract}

Following Proposition 4.1 of the foundation paper, every finite word $W$ (length $\ell=|W|$, $k=|W|_1$ ones) has an explicit periodic-value formula
\[
\Phi(W^\infty) = \frac{c_W}{2^\ell - 3^k}
\]
for an explicit integer $c_W$ (\repo{THEOREM\_STATUS.md}, citing the foundation paper's Prop.\ 4.1).

Suppose, for contradiction, $\Phi(s)=u/v\in\Q\cap\Z_2$ in lowest terms. For any finite word $W_n$ with $s\ne W_n^\infty$, put $\ell_n=|W_n|$, $k_n=|W_n|_1$, $L_n=\operatorname{lcp}(s,W_n^\infty)$, $\delta_n = 2^{\ell_n}-3^{k_n}$, and define the integer
\[
M_n := u\,\delta_n - v\,c_{W_n}.
\]
By the 2-adic isometry, $v_2(\Phi(s)-\Phi(W_n^\infty))\ge L_n$, which translates into a 2-adic divisibility lower bound on $M_n$; on the other hand $|M_n|$ is bounded above in terms of $|u|,v,\delta_n,c_{W_n}$ --- an Archimedean (ordinary absolute-value) bound. When the 2-adic lower bound forces $M_n$ to be divisible by more than its Archimedean bound allows, $M_n=0$ is forced, and (with $s\ne W_n^\infty$) that is eventually impossible. This comparison is exactly what the repository's \emph{surplus} quantity measures.

\begin{definition}[Surplus, canonical form \cite{Phase2Surplus}]
\[
S_s(W) := \operatorname{lcp}(s,W^\infty) - \log_2\bigl(|2^{|W|}-3^{|W|_1}| + c_W\bigr).
\]
\end{definition}
This is the repository's own proof-native quantity, fixed from Phase 1 onward and not altered by any later phase: \emph{``the proof-native quantity, not a convenient after-the-fact packaging''} (\repo{theorems/phase2/SURPLUS\_THEOREM.md}).

\begin{theorem}[Abstract 2-adic periodic-approximation criterion \cite{Phase2Abstract}]
\label{thm:abstract}
Let $s\in\{0,1\}^\N$ and $(W_n)_{n\ge1}$ finite binary words with $s\ne W_n^\infty$ for all $n$. If
\[
\limsup_{n\to\infty} S_s(W_n) = +\infty,
\]
then $\Phi(s)\notin\Q$.
\end{theorem}
This is stated and proved with no Sturmian or Rote terminology at all --- it is the Sturmian-free abstraction every later result in this note specializes.

\begin{theorem}[Dual necessity, \cite{Phase2Surplus}]
If $\Phi(s)=u/v\in\Q$ in lowest terms, height $H=\max(|u|,v)$, then for every finite $W$ with $s\ne W^\infty$,
\[
S_s(W) \le \log_2 H.
\]
\end{theorem}
Together: \emph{surplus $\to+\infty$ along some sequence of periodic approximants forces irrationality; rationality forces surplus to stay uniformly bounded.} Everything that follows is the construction of periodic approximants along which surplus is provably unbounded, for successively broader symbolic classes.

\section{Discrepancy as the Height Bridge}
\label{sec:discrepancy}

Repetition length $L_n$ alone is not enough: the surplus definition also involves $c_W$, and the periodic-value formula gives no a priori control on how large $c_W$ can be relative to $2^\ell-3^k$. The quantity that closes this gap is prefix \emph{discrepancy}:
\[
D(W) := \max_{0\le i\le \ell} \left| k_i(W) - \frac{ik}{\ell}\right|,
\]
where $k_i(W)$ counts ones in the length-$i$ prefix of $W$.

\begin{theorem}[Discrepancy-height bound \cite{Phase1Discrepancy}]
For every finite word $W$, $\ell=|W|$, $k$ ones, $D=D(W)$:
\[
c_W \;\le\; \ell\cdot 3^{\lceil D\rceil}\cdot\max(2^\ell,3^k).
\]
\end{theorem}
(Proved, and exhaustively verified by direct computation for $\ell\le 16$.) This single inequality is what converts symbolic regularity (a bound on how unevenly $1$s are distributed in $W$) into an arithmetic height estimate the surplus computation can use. The relevant hierarchy of discrepancy regimes, from weakest to strongest sufficient hypothesis actually used somewhere in this program, is:
\begin{itemize}
\item $D(W)=O(1)$ --- e.g., single-block or highly regular words; strongest, rarely needed.
\item $D(W)=O(\log \ell)$ --- achieved for rotations by any slope with \emph{bounded partial quotients}, via the classical three-distance theorem; this is the regime used throughout \S\ref{sec:oddweight}--\S\ref{sec:boundedtype}.
\item $D(W)=o(\ell)$ --- the general sufficient hypothesis for the abstract sublinear-discrepancy theorem below; strictly weaker, but gives up the explicit logarithmic rate.
\end{itemize}
Sublinear discrepancy is exactly what is needed for the height cost to be dominated by a linear repetition gain (\S\ref{sec:repheight}); the sharper $O(\log\ell)$ rate is what the bounded-partial-quotient hypothesis buys on top of that.

\section{Repetition--Height Competition}
\label{sec:repheight}

\begin{theorem}[Sublinear-discrepancy repetition criterion \cite{Phase1Discrepancy}]
Suppose $s$ begins in initial powers $W_n^{r_n}$, $\ell_n=|W_n|\to\infty$, densities $\gamma_n=k_n/\ell_n\to\gamma_\infty$, and $D(W_n)=o(\ell_n)$. If
\[
\liminf_n r_n > \max(1,\gamma_\infty\log_2 3) \quad\text{strictly},
\]
then $S_s(W_n)\to+\infty$, hence $\Phi(s)\notin\Q$.
\end{theorem}
Taking the supremum over $\gamma\in(0,1)$ of $\max(1,\gamma\log_2 3)$ gives $\log_2 3=1.58496\ldots$, so $r>\log_2 3$ suffices \emph{uniformly in the density}. Because $\log_2 3<2$, an initial \emph{square} ($r=2$) always clears this bar, with an explicit margin $2-\log_2 3=0.41503\ldots$ --- which is exactly why squares are the workhorse repetition structure throughout this program:

\begin{theorem}[Squares suffice unconditionally \cite{Phase1Repetition}]
If $s$ begins in arbitrarily long initial squares $W_nW_n$ --- with \emph{no} hypothesis on discrepancy, balance, or internal structure of the roots $W_n$ --- then $\Phi(s)\notin\Q$.
\end{theorem}

The later, more flexible sufficient condition used once a fixed margin is no longer available (\S\ref{sec:ostrowski}) does not require a fixed constant margin at all; it only requires the repetition \emph{excess} above the threshold to dominate the discrepancy-plus-logarithmic height cost, i.e.\ excess $\cdot\, \ell \gg \log \ell$ rather than excess $\ge$ const.

\section{The Sturmian Foundation}
\label{sec:sturmian}

The foundation paper \cite{DeJesus2026Sturmian} proves, by \emph{two independent mechanisms}, that every Sturmian parity word has irrational $\Phi$-value:

\begin{enumerate}
\item \textbf{Counting.} Every Sturmian word has $p(L)=L+1$, far below the Dubickas threshold, so Route A (\S\ref{sec:routeA}) applies directly.
\item \textbf{Liouville / initial-square route.} Sturmian words have arbitrarily long initial squares (Theorem 10.2 of \cite{DeJesus2026Sturmian}, citing Allouche, Davison, Queffélec, and Zamboni \cite{ADQZ2001}), which triggers the repetition--height mechanism of \S\ref{sec:repheight} directly (Theorem 10.1 of \cite{DeJesus2026Sturmian}). The foundation paper also gives a second, independent and effective proof, for the specific case López and Stoll's own paper \cite{LopezStoll2009} had identified as open (their critical-density Sturmian word at slope $\beta=\ln2/\ln3$): this is then generalized to every irrational slope, and a further argument removes the restriction to intercept $0$ via an initial square of the word itself.
\end{enumerate}

Both mechanisms are cited, not re-derived, in this note; nothing in this program alters or extends the Sturmian case itself. What matters for everything that follows is where the foundation paper's own abstract explicitly places the frontier: \emph{``Counting reaches every complexity below $1.70951\ldots L$; the frontier is complexity $2L$ (Rote words, rotation codings), open for both routes.''} This is the exact motivation for the program consolidated here.

\section{Crossing the Counting Frontier: Complementary Symmetric Rote Sequences}
\label{sec:rote}

A \emph{complementary symmetric (CS) Rote sequence} $v\in\{0,1\}^\N$ is one whose first-difference (XOR) transform is Sturmian:
\[
v \text{ is CS Rote} \iff u:=S(v) \text{ is Sturmian}, \qquad u_i = v_i\oplus v_{i+1}\ \ (i\ge0),
\]
with inverse recovery $v_{n+1}=v_n\oplus u_n$ \cite{Phase2Transfer}. Such sequences have factor complexity exactly $p(n)=2n$, standard in the combinatorics-on-words literature on this class (background on the surrounding Sturmian/complementary-word theory: \cite{Lothaire2002}; see the bibliography note on primary Rote-sequence sourcing) --- past the Dubickas threshold, so Route A never applies to them; every irrationality result about CS Rote sequences in this note is a Route B result.

\begin{theorem}[Exact transfer, with a $+1$ bonus \cite{Phase2Transfer}]
\label{thm:transfer}
Let $L:=\operatorname{lcp}(u,W^\infty)$ for a Sturmian $u=S(v)$ and finite word $W$, and let $L_v:=\operatorname{lcp}(v,\mathrm{target}^\infty)$ where $\mathrm{target}=V$ (if $W=V$ has even weight) or $\mathrm{target}=R:=V\bar V$ (if $V$ has odd weight, $\bar V$ its bitwise complement). Then
\[
L_v = L+1, \quad\text{exactly, not merely } \ge L \text{ or } \approx L.
\]
\end{theorem}
Verified by induction (both directions of the inequality) and computationally over $6$ initial cases in Phase 2; re-derived completely independently in Phase 5 and cross-checked exhaustively over every binary root of length $2$ to $10$ and every mismatch position from $\ell$ to $4\ell$ --- $114{,}680$ cases, $0$ failures \cite{Phase5Transfer}. The theorem is intercept-agnostic: it holds verbatim regardless of the intercept of the underlying Sturmian word.

\section{The Odd-Weight Mechanism}
\label{sec:oddweight}

This is the structural discovery that makes the whole bounded-type theorem (\S\ref{sec:boundedtype}) possible, and it deserves to be isolated on its own.

When the Sturmian root $W$ has \emph{odd} weight, Theorem \ref{thm:transfer}'s target is $R=V\bar V$, not $V$ itself.

\begin{theorem}[Density of the doubled root, exactly $1/2$ \cite{Phase4Required}]
For every finite word $V$, $R:=V\bar V$ satisfies $|R|=2|V|$ and $|R|_1=|V|$, hence
\[
\text{density}(R) = \frac{|R|_1}{|R|} = \frac12, \quad\text{exactly, unconditionally.}
\]
\end{theorem}
\begin{proof}[Proof (as recorded)]
$|R|_1 = |V|_1 + |\bar V|_1 = |V|_1 + (|V|-|V|_1) = |V| = \ell$; and $|R|=2\ell$; so density $=\ell/(2\ell)=1/2$. $\blacksquare$
\end{proof}

Arithmetically, this matters because $3^{|R|_1}=3^{|R|/2}$ grows strictly more slowly than $2^{|R|}$ (since $\sqrt3<2$), so the dominant exponential height scale in the periodic-value denominator $2^{|R|}-3^{|R|_1}$ is unambiguously $2^{|R|}$, regardless of the fine structure of $V$. Plugging $\gamma=1/2$ into the repetition--height criterion of \S\ref{sec:repheight} collapses the required threshold on the Sturmian side from an earlier, over-conservative estimate of $2\log_2 3\approx3.170$ down to the true requirement,
\[
r_{\text{required}}(\gamma) = 2, \qquad\text{uniformly, for both the even- and odd-weight transfer routes.}
\]
This is established in \cite{Phase4Required}. This single fact replaced two working hypotheses this program had earlier formed and found false --- a universal even-weight witness hypothesis, and a silver-ratio extremality claim --- both recorded in full in \S\ref{sec:selfcorrection}.

\section{Self-Correction History}
\label{sec:selfcorrection}

This program is falsification-first throughout: every phase actively attacked its own prior claims before extending them. Four corrections are recorded here in full, not hidden, because they show how the final theorem (\S\ref{sec:boundedtype}) was actually reached, and because a future researcher extending this work needs to know which paths were tried and ruled out.

\paragraph{A. Universal even-weight witness hypothesis --- rejected.} An early working assumption was that the Rote-transfer route could always be arranged to use an even-weight root $V$ (avoiding the doubled root $R$ altogether). A finite $8$-state automaton analysis of the weight-parity sequence found exactly one ``bad'' cycle sustaining all-odd weight forever, and an explicit quadratic-irrational counterexample slope was found and verified to $28$ convergents: $\gamma=[0;1,1,1,\overline{2}]$ has weight parity $0,0,1,1,1,1,1,\ldots$ --- all odd (weight $1$) from index $2$ onward, forever \cite{Phase3Automaton}. The universal claim is false. The bridge still succeeds on this slope because the odd-weight transfer route (\S\ref{sec:oddweight}) handles exactly this case.

\paragraph{B. Silver-ratio extremality --- rejected, explicit counterexample.} A subsequent working hypothesis, formed from testing only four continued-fraction tails, was that among all-even-tail slopes the silver ratio ($\overline{2}$ tail) is \emph{extremal}, with minimal achievable exponent $2+\sqrt2\approx3.4142$. An adversarial search found $\gamma=[0;1,4,1,5,1,\overline{6,2}]$ with $\text{ice}(u)\approx3.1547 < 2+\sqrt2$ \cite{Phase4Silver}, disproving the hypothesis outright. The constant $2+\sqrt2$ itself is not wrong --- it is the exact value for the pure silver-ratio tail \cite{Phase4SilverAudit} --- but it was the wrong threshold to be solving for: once the density-$1/2$ fact (\S\ref{sec:oddweight}) is known, the actual required threshold is $r>2$, and $2+\sqrt2\gg2$ was never close to the real boundary.

\paragraph{C. Seed-independence via complement-invariance of $c_W$ --- rejected, repaired.} The claim that both CS Rote seeds ($v_0=0$ or $v_0=1$) give identical surplus was correct, but its stated justification (``complementing uniformly complements $R$, so $c_W$ is unaffected'') was not: direct testing found $c_W\ne c_{\bar W}$ in $8$ of $8$ random odd-weight cases, often by an order of magnitude \cite{Phase5Hostile}. The repaired reason: every verification script's height bound $F_{\text{bound}}(W)$ depends on the root only through $(\ell,k,D(\text{root}))$, never through $c_W$ directly, and $D(V\bar V)=D(\bar V V)$ exactly, checked in $15$ of $15$ random cases. The \emph{conclusion} survives untouched; only its justification needed repair --- exactly the kind of error a hostile self-audit is designed to catch.

\paragraph{D. Phase 6 small-$k$ boundary correction.} The load-bearing inequality $\text{term2}(k)>2+1/(A+1)$ (\S\ref{sec:boundedtype}) was stated to hold strictly for every $k$. Adversarial testing against genuinely non-eventually-periodic bounded-type controls (Thue--Morse-coded, Fibonacci-word-coded, and two independent random $\{1,2\}$ sequences) found one case of non-strict \emph{equality} rather than a violation: in the Fibonacci-coded control, $\text{term2}(2)$ equals $2+1/(A+1)=7/3$ exactly, at $k=2$, the only such case out of $232$ checks across all four sequences \cite{Phase6Arrow}. This does not weaken the theorem, because the abstract criterion (Theorem \ref{thm:abstract}) only requires $\limsup$ unboundedness along \emph{some} growing sequence of witnesses, not strict inequality at literally every finite $k$.

None of these four corrections changed the truth of the final headline theorem; each changed (and, in cases B and C, simplified) the reason it is true.

\section{The Bounded-Type CS Rote Theorem}
\label{sec:boundedtype}

\begin{theorem}[Headline result of the program \cite{Phase6Bounded}]
\label{thm:boundedtype}
Let $\gamma\in(0,1)$ be irrational with bounded partial quotients ($A:=\sup_n a_n<\infty$). Let $u$ be the intercept-$0$ lower mechanical word of $\gamma$. Let $v$ be either complementary symmetric Rote sequence with $S(v)=u$ (either seed), or any shift $\sigma^j(v)$. Then
\[
\Phi(v)\notin\Q.
\]
\end{theorem}

This supersedes an earlier theorem of the same program, proved first for the strictly narrower class of \emph{quadratic}-irrational slopes \cite{Phase4Quadratic}, \emph{without weakening any of its quantifiers} --- both seeds, every shift, still covered. The upgrade matters because quadratic irrationals are countable (each is a root of an integer quadratic, and there are only countably many such polynomials), while bounded-partial-quotient irrationals are uncountable: for each bound $A$, the set of irrationals whose partial quotients never exceed $A$ contains a Cantor-type subset, and the union over all $A$ is an uncountable set.

\begin{remark}[What ``uncountable'' does and does not say here]
Uncountability is a statement about cardinality only. The set of bounded-type irrationals has Lebesgue measure zero: a Lebesgue-almost-every real number has \emph{unbounded} partial quotients. Theorem \ref{thm:boundedtype} says nothing about a randomly chosen real slope, and nothing about density among all irrationals --- only that the class it covers, while structurally constrained, is not merely a countable list of algebraic special cases.
\end{remark}

\section{Proof Architecture of the Bounded-Type Theorem}
\label{sec:architecture}

The shortest audited chain, reconstructed in Phase 6 without any periodicity hypothesis and cross-checked against four genuinely non-eventually-periodic bounded-type controls \cite{Phase6Arrow}, adapting the original Phase 4 chain \cite{Phase4Chain}:

\begin{center}
\begin{tabular}{ll}
\toprule
Step & Cites \\
\midrule
bounded partial quotients, $A=\sup_n a_n$ & hypothesis \\
$\downarrow$ Sturmian words have arbitrarily long initial squares & Thm 10.2 \cite{DeJesus2026Sturmian}, \cite{ADQZ2001} \\
$\downarrow$ BHZ initial-power margin, $\text{term2}(k)>2+\tfrac1{A+1}$ (every $k$, up to \S\ref{sec:selfcorrection}.D) & \cite{BHZ2006}, Cor.\ 3.5 \\
$\downarrow$ arbitrarily long usable witnesses & \S\ref{sec:repheight} \\
$\downarrow$ exact XOR transfer, $L_v=L+1$ & Thm \ref{thm:transfer} \cite{Phase2Transfer,Phase5Transfer} \\
$\downarrow$ even/odd branch split & \S\ref{sec:rote} \\
$\downarrow$ odd branch: doubled root density exactly $1/2$ & \S\ref{sec:oddweight} \cite{Phase4Required} \\
$\downarrow$ actual-root discrepancy control, $O(\log\ell)$ & \cite{Phase2Discrepancy}, repaired \cite{Phase5Hostile} \\
$\downarrow$ sublinear height penalty & \S\ref{sec:discrepancy} \\
$\downarrow$ positive, unbounded surplus & \S\ref{sec:abstract} \\
$\downarrow$ abstract irrationality criterion (Thm \ref{thm:abstract}) & \cite{Phase2Abstract} \\
$\Rightarrow$ $\Phi(v)\notin\Q$ & Thm \ref{thm:boundedtype} \\
\bottomrule
\end{tabular}
\end{center}

Notably absent from this chain: the weight-parity automaton of \S\ref{sec:selfcorrection}.A, and the silver-ratio constant $2+\sqrt2$ of \S\ref{sec:selfcorrection}.B --- both were real, correctly-derived facts about this problem, and neither is needed by the final proof.

\section{General Intercepts}
\label{sec:generalintercept}

Everything in \S\ref{sec:oddweight}--\S\ref{sec:architecture} is stated for the \emph{intercept-$0$} representative of a slope $\gamma$: the lower mechanical word $u_i=\lfloor(i+1)\gamma\rfloor-\lfloor i\gamma\rfloor$. Removing this restriction --- treating CS Rote sequences built on \emph{any} intercept $\rho$ of the same slope --- is where the program currently stops short of a full theorem.

An arbitrary intercept is encoded by an \emph{Ostrowski digit sequence} $(c_k)_{k\ge1}$, subject to the admissibility condition \cite{BHZGeneralMap}:
\[
a_k \ge c_k \ge 0 \text{ for all } k,\quad a_k\ge1 \text{ for } k\ge2,\quad \text{and if } c_k=a_k \text{ then } c_{k-1}=0.
\]
The characteristic (intercept-$0$) sequence is $c_k\equiv0$. Where the partial quotients $(a_k)$ fix the canonical Sturmian \emph{scale hierarchy} (the sequence of convergent denominators $q_k$), the digits $(c_k)$ determine \emph{where inside that hierarchy} a particular intercept's word sits.

At intercept $0$, Berthé--Holton--Zamboni's Corollary 3.5 collapses to the clean closed form used in \S\ref{sec:architecture}. In general, it does not collapse:
\[
x'(k) = \frac{\sum_{j=1}^{k+1}(a_j-c_j)q_{j-1}}{q_k}, \qquad
y(k) = 1 + \frac{\sum_{j=1}^{k}(a_j-c_j)q_{j-1}}{q_k-c_kq_{k-1}},
\]
\[
\text{ice}(\omega)=\limsup_k \max(x'(k),y(k)).
\]
Both terms now depend on the entire digit history $(c_j)$, not just on $(a_j)$, and the simplification that made $\text{term2}(k)>2+1/(A+1)$ hold cleanly at intercept $0$ does not transfer automatically.

This is precisely and honestly where the program's current record stops:

\begin{conjecture}[The one remaining lemma \cite{Phase7Verdict}]
\label{conj:phase7}
For any admissible Ostrowski digit sequence $(c_k)$ of a bounded-partial-quotient slope $\gamma$ ($A=\sup_n a_n<\infty$),
\[
\limsup_k \max\bigl(x'(k),y(k)\bigr) > 2,
\]
with a margin depending only on $A$.
\end{conjecture}

This is classified \emph{conjectural} in the repository record, with strong, broad computational evidence ($36$ of $36$ tested configurations, including the family identified as the empirical worst case, across a range of $A$) and \emph{no completed algebraic proof} \cite{Phase7Sufficient}. It is \textbf{not} claimed proved here or anywhere in the underlying record. If Conjecture \ref{conj:phase7} is established, Theorem \ref{thm:boundedtype} would extend verbatim to every intercept of every bounded-type slope; until then, that extension remains open.

The repository's own Phase 7 verdict summarizes the state of this obstruction precisely: \emph{``The obstruction \ldots is precisely located, its downstream consequences fully traced, and its practical severity found to be much smaller than initially suspected --- no computational evidence of failure anywhere \ldots''} \cite{Phase7Verdict}.

\section{Ostrowski Repetition Structure}
\label{sec:ostrowski}

A subtlety worth making explicit for anyone continuing this work: \textbf{low factor complexity is not the same thing as strong initial repetition.} A CS Rote sequence may remain globally $p(n)=2n$ while its intercept changes \emph{which} large canonical blocks (Ostrowski/continuant-scale factors) happen to align at the origin --- which is exactly what determines how long an initial repetition is actually available to feed the bridge.

The scale-dependent quantity that tracks this is an \emph{excess}:
\[
E_k := \max(r'_k, r_k) - 2,
\]
where $r'_k,r_k$ are the two branch-specific repetition exponents at scale $k$ \cite{Phase7Surplus}. The corresponding surplus decomposition is
\[
S_k \;\ge\; A\cdot E_k \;-\; C_A\cdot D_k \;-\; O_A(\log \ell_k)
\]
for constants depending only on $A$ \cite{Phase7Surplus}. The flexible sufficient condition this leads to --- used in place of a fixed margin once one is no longer cleanly available --- is
\[
E_k \gg \frac{\log \ell_k}{\ell_k}.
\]
This is the repository's own formulation of the sufficient condition \cite{Phase7Flexible}.
Every structured intercept family actually tested lands comfortably inside this condition, with no case found close to violating it \cite{Phase7Structured}. This excess-versus-log-height competition, not a fixed numeric threshold, is the current frontier condition for the general-intercept problem, and is the natural object for a future proof of Conjecture \ref{conj:phase7} to control directly.

\section{Interpretation: Symbolic, Statistical, and Arithmetic Space}
\label{sec:interpretation}

This section is conceptual, not a theorem statement, and should be read as orientation for a reader extending this work rather than as a claim.

It is useful to separate three levels at which a parity word can be described:

\begin{description}[leftmargin=2em]
\item[Symbolic space.] The actual word: its literal order, block structure, and factor-by-factor arrangement.
\item[Statistical projection.] Aggregated summaries of the symbolic object --- density $\gamma=k/\ell$, discrepancy $D(W)$, block frequencies, factor complexity. Many distinct symbolic objects can share the same statistical projection.
\item[Arithmetic admissibility.] Which symbolic objects, specifically, correspond to a rational $\Phi$-value. This is a pointwise, not a statistical, property.
\end{description}

The relationship is a \emph{fiber}: symbolic object $\to$ statistical state is generally many-to-one. A crucial consequence, easy to state and easy to get wrong in either direction: \textbf{statistical sparsity of a class (low complexity, small discrepancy, controlled density) does not by itself imply arithmetic impossibility of rational $\Phi$-values within it}, and conversely, statistical richness does not by itself force irrationality. The bridge proved in this note works precisely because it does \emph{not} stop at a statistical argument: it lifts symbolic and statistical information (repetition length, discrepancy) back into a pointwise arithmetic contradiction (Theorem \ref{thm:abstract}'s exact integer divisibility argument) for the \emph{specific} word $s$ under study, not for a statistical ensemble containing it. This is why the method scales from Sturmian words to CS Rote sequences without a change in kind --- and also why it stalls exactly where the statistical control (discrepancy, excess) it depends on stalls (\S\ref{sec:generalintercept}).

\section{The Drift/Resonance Coordinate}
\label{sec:drift}

For a finite word $W$ ($\ell=|W|$, $k=|W|_1$), define
\[
R(W) := \ell - k\log_2 3.
\]
Then $2^\ell - 3^k = 2^\ell\bigl(1-2^{-R(W)}\bigr)$, and near $R=0$ (i.e.\ near the ``resonance'' $\ell\approx k\log_2 3$), $|2^\ell-3^k| \approx 2^\ell \ln2\,|R(W)|$.

This coordinate is explanatory, not structural: it was investigated directly in Phase 7 as a candidate new invariant, and the record is explicit that it is not one. \emph{``No new information is encoded by $R$ that $\gamma$ did not already carry''} \cite{Phase7Drift}; \emph{``Resonance did not materially help or hurt in any tested case''} \cite{Phase7DriftDigits}. $R(W)$ is a reparameterization of the density margin $\gamma\log_2 3$ already used throughout \S\ref{sec:repheight}--\S\ref{sec:oddweight}, useful for visualizing how close a given witness sits to the arithmetic boundary, but supplying no additional proof-theoretic leverage found anywhere in this program.

\section{The Symmetry-Tradeoff Probe (Exploratory, Inconclusive)}
\label{sec:symmetryprobe}

\textbf{This section reports an exploratory investigation. Nothing in it is a theorem, and nothing in it should be read as evidence for or against any conjecture stated elsewhere in this note.} It is included so a future researcher does not repeat the same probe without knowing what was already tried and why it did not resolve anything.

The idea investigated: does \emph{approach to symmetry} in the Ostrowski digit sequence (roughly, how often the maximal admissible continuation is available) trade off against arithmetic \emph{resonance} (how close $R(W)$, \S\ref{sec:drift}, sits to $0$)? Two experiments were run.

\textbf{Experiment 1} correlated a normalized arithmetic-slack quantity $A(W):=|1-2^{-R(W)}|$ against the exact surplus $S(W)$ and against the achieved repetition ratio $L/\ell$, using $74$ genuine witnesses from the general-intercept search of \S\ref{sec:generalintercept}. Result: $A(W)$ vs.\ $S(W)$, $r\approx0.208$; $A(W)$ vs.\ $L/\ell$, $r\approx-0.125$. Both fall below the probe's own $0.3$ threshold for ``worth a theorem-seeking follow-up,'' and $n=74$ is too small for $r=0.208$ to clear the standard significance threshold for that sample size. Separately, the first correlation is only weak evidence in any case: the surplus machinery already leans on the same $\ell-k\log_2 3$ resonance quantity that defines $A(W)$, so a correlation there may be partly definitional rather than an independent discovery.

\textbf{Experiment 2}, testing a genuinely independent quantity (a windowed frequency of admissible-digit freedom, $N_{\text{sym}}$, against $A(W)$), did not run to completion in its original form: it materialized an explicit substitution-image word whose length grows exponentially in the search depth $k$ (for slope $A=6$, a projected $\sim9.9\times10^{14}$ symbols by $k=19$), exhausting memory. This was diagnosed and repaired: the word's length and weight can be obtained exactly, without ever materializing it, via a $2\times2$ integer transfer-matrix representation of the same substitutions ($\tau_0,\tau_1$; see \S\ref{sec:transfermatrix}), cross-checked against brute-force construction in $150$ of $150$ small cases with exact agreement. With this fix, Experiment 2 completed on $724$ points: $A(W)$ vs.\ $N_{\text{sym}}$, $r\approx-0.0058$; on a log--log scale, $r\approx-0.0319$. \textbf{These numbers must not be read at face value}: $596$ of the $724$ points ($82\%$) had $A(W)$ pinned to exactly $1.0$ by the probe's own floating-point underflow clamp ($|R(W)|>60$), an artifact of the diagnostic's float arithmetic, not a genuine measurement in the near-resonance regime where any real effect would have to live.

\begin{quote}
\textbf{Verdict: INCONCLUSIVE.} Neither experiment produces a correlation free of confounding structure (definitional entanglement in Experiment 1; a numerical clamping artifact dominating Experiment 2). This is consistent with, and independently corroborates, the Phase 7 drift-coordinate finding of \S\ref{sec:drift} that resonance carries no separately exploitable signal in every case tested so far --- but it establishes nothing new, and \emph{the symbolic-symmetry/arithmetic-freedom tradeoff hypothesis is neither supported nor contradicted} by this probe.
\end{quote}

A methodologically sound re-run would need $A(W)$ computed at sufficient precision to avoid the clamp, restricted to (or stratified by) genuine near-resonance witnesses, before any correlation from it could be informative.

\section{A Transfer-Matrix Computational Tool}
\label{sec:transfermatrix}

\textbf{This section records a computational/structural tool, not a new irrationality theorem.} It is preserved here because it is broadly useful and easy to lose track of.

The Sturmian substitutions used throughout this program (e.g., in constructing explicit BHZ witness words, \S\ref{sec:generalintercept}) are
\[
\tau_0: 0\mapsto 0,\ 1\mapsto 01; \qquad \tau_1: 0\mapsto 10,\ 1\mapsto 1.
\]
Applying either substitution repeatedly can make the resulting word's length grow exponentially in the number of applications --- and that length is all that is needed for most downstream computations (arithmetic slack $A(W)$, weight $\kappa$, density). It need never be materialized. Tracking only $(n_0,n_1)$, the counts of $0$s and $1$s, the two substitutions act linearly:
\[
\tau_0: (n_0,n_1)\mapsto(n_0+n_1,\ n_1), \qquad
\tau_1: (n_0,n_1)\mapsto(n_0,\ n_0+n_1),
\]
i.e.\ by the integer matrices $M_0=\begin{psmallmatrix}1&1\\0&1\end{psmallmatrix}$, $M_1=\begin{psmallmatrix}1&0\\1&1\end{psmallmatrix}$ (both determinant $1$; their products are exactly the continuant recurrence generating the convergent denominators $q_k$). Applying $\tau_i$ exactly $e$ times in a row is $M_i^e$; composing across an entire substitution schedule of length $k$ is $k$ integer matrix--vector multiplications, giving the word's exact length $\ell=n_0+n_1$ and weight $\kappa=n_1$ in $O(k)$ exact-integer steps regardless of how astronomically large $\ell$ becomes --- verified against brute-force construction with $150/150$ exact agreement wherever brute force was still feasible to run.

A related, separate observation: counting the number of admissible Ostrowski digit sequences $(c_k)$ (rather than computing one fixed sequence's word) is itself governed by a minimal two-state automaton, not by any larger state. The admissibility rule (\S\ref{sec:generalintercept}) depends only on whether the immediately preceding digit was exactly $0$: from that state, the next digit has $a_k+1$ admissible choices; from any nonzero-previous state, only $a_k$. No memory of the actual previous digit value, nor of anything beyond the current index into $(a_k)$, is needed. This directly mirrors the admissible-sequence generator already used throughout the repository's scripts, and is recorded here as a structural fact about that generator, not as a separately theorem-proving result.

\section{What Is Still Open}
\label{sec:open}

\begin{description}[leftmargin=3em,style=nextline]
\item[Open 1 --- General intercept margin.] Prove or refute Conjecture \ref{conj:phase7}: the BHZ/Ostrowski margin lemma for arbitrary admissible digit sequences on bounded-type slopes (\S\ref{sec:generalintercept}).
\item[Open 2 --- Return/exceptional intercept set.] The general-intercept BHZ formula identifies a countable ``exceptional'' or return-intercept set ($\rho=\{j\gamma\}$) with a different, generally larger, margin structure \cite{BHZGeneralMap}. Classify this set precisely and determine whether every member still supplies unbounded surplus.
\item[Open 3 --- Intrinsic-square bypass.] Investigate whether arbitrary-intercept intrinsic squares of the word itself (in the spirit of the foundation paper's own initial-square route, \S\ref{sec:sturmian}) can bypass the full BHZ/Ostrowski machinery entirely, rather than requiring Conjecture \ref{conj:phase7}. Preliminary evidence: bare squares ($r=2$) are already known to be insufficient on the odd branch even at intercept $0$; strict excess $r>2$ is required there too \cite{Phase7Intrinsic}, so a bypass, if one exists, is not merely ``find any square.''
\item[Open 4 --- General CS Rote.] Remove the bounded-partial-quotient hypothesis from Theorem \ref{thm:boundedtype} if possible. Flagged in the repository record as genuinely out of reach of the discrepancy mechanism used here, since the $O(\log\ell)$ discrepancy bound for rotations is a bounded-partial-quotient-specific consequence of the three-distance theorem.
\item[Open 5 --- Higher complexity.] Only once the $p(n)=2n$ rung is fully closed (Open 1--4): which symbolic classes beyond Rote sequences --- higher, structured complexity classes; general two-interval rotation codings, for which the repetition-existence half of the argument is itself still open \cite{Phase2Rotation}; general morphic words --- can this bridge reach next?
\item[Open 6 --- The full Periodicity Conjecture.] What mechanism, if any, could cover aperiodic words that admit no useful initial repetition structure at all? The Thue--Morse word has no square whatsoever below half-length $2048$ \cite{Phase1Gap}, making every method in this note --- indeed, every purely repetition-based method --- structurally inapplicable to it by construction. This is not a gap to be closed by extending the present technique further; it is evidence that a fundamentally different mechanism would be needed for the conjecture in full generality.
\end{description}

\section{Future-Work Roadmap}
\label{sec:roadmap}

The roadmap below is organized directly around the six open problems of \S\ref{sec:open} and the repository's own stated priorities.

\begin{description}[leftmargin=3em,style=nextline]
\item[Near term (directly attackable with current machinery).] Attempt a full algebraic proof of Conjecture \ref{conj:phase7} (Open 1), most promisingly by pushing the ``exceptional-orbit reduction'' already used to narrow the genuinely open cases (\repo{theorems/phase7/PHASE7\_VERDICT.md}) the rest of the way, rather than by further enlarging the computational search that has already found $36/36$ supporting configurations and no counterexample. In parallel, classify the countable return-intercept set of Open 2 explicitly enough to check it case-by-case against whatever form Conjecture \ref{conj:phase7}'s proof takes.
\item[Medium term.] Investigate the intrinsic-square bypass of Open 3 as an independent, possibly simpler, route to the same general-intercept conclusion --- worth pursuing in parallel with, not only as a fallback from, Open 1. If it succeeds, re-derive Theorem \ref{thm:boundedtype}'s general-intercept extension from it directly and check whether the two routes agree everywhere they overlap.
\item[Longer term.] Once intercepts are fully closed, decide the shape of Open 4 (general CS Rote, unbounded partial quotients) before attempting Open 5: the discrepancy mechanism this program uses is specifically a bounded-partial-quotient consequence, so extending past it likely requires either a genuinely different discrepancy estimate for unbounded-type rotations, or an argument that does not route through discrepancy at all.
\item[Structural.] Any attempt at Open 6 (the full conjecture) needs a mechanism that does not depend on initial repetition existing at all, given the Thue--Morse-type structural ceiling recorded in \S\ref{sec:open}. This program takes no position on what such a mechanism could look like; it records only that repetition-based methods, however far extended, cannot be it alone.
\item[Process.] Continue the falsification-first discipline of \S\ref{sec:selfcorrection}: every extension of this program should include an adversarial audit pass against genuinely non-periodic, non-quadratic controls before any theorem upgrade is accepted, exactly as Phases 5 and 6 did for the bounded-type theorem.
\end{description}

\section*{Provenance and Verification}
\addcontentsline{toc}{section}{Provenance and Verification}

This note was consolidated from the \repo{periodicity-conjecture-bridge} repository at commit \texttt{a4ac46aebf17fb2ce14cea50f5ce60a58704da15} (``Phase 7: map general-intercept CS Rote obstruction''), one commit ahead of \texttt{origin/main} and not pushed during this consolidation. Every theorem statement, definition, and quoted result above was extracted directly from the corresponding file under \repo{theorems/phase\{1..7\}/} and cross-checked against \repo{THEOREM\_STATUS.md} and \repo{LADDER.md}; none was reconstructed from memory. SHA-256 checksums of every file under \repo{theorems/}, \repo{scripts/}, \repo{data/}, and \repo{reference/} at the time of writing are recorded in \repo{technical-note/PRESERVATION\allowbreak\_RECORD.md}. A full statement-by-statement classification (known/cited, proved in foundation paper, proved in this program, computationally verified, counterexample, rejected hypothesis, open, or exploratory only) is recorded in \repo{technical-note/SOURCE\_OF\_TRUTH.md}. No Phase 1--7 research file was modified in the preparation of this note.

This note does not include a Zenodo DOI. Once a DOI is minted for this deposition, it should be added to \repo{CITATION.cff} and to this document's front matter; \repo{technical-note/ZENODO\_PREP.md} contains ready-to-use snippets for both, prepared but not applied.

\section*{Acknowledgments}
\addcontentsline{toc}{section}{Acknowledgments}
AI assistance (Anthropic's Claude) was used for repository archaeology, adversarial diagnostic work (including diagnosing and repairing the Experiment 2 performance bug in \S\ref{sec:symmetryprobe}), literature cross-checking against the foundation paper's own bibliography, and drafting of this note, consistent with the acknowledgment practice of the foundation paper \cite{DeJesus2026Sturmian}. All mathematical claims, their framing, and any errors are the author's.

\begin{thebibliography}{99}

\bibitem{DeJesus2026Sturmian} Elias De Jesús, \emph{The 3x+1 Conjugacy Map Sends Every Sturmian Word to an Irrational 2-Adic Integer}, unpublished manuscript, September 2026. (Foundation paper for this research program; companion work by the present author.)

\bibitem{ADQZ2001} Jean-Paul Allouche, J.\ L.\ Davison, Michel Queffélec, and Luca Q.\ Zamboni, \emph{Transcendence of Sturmian or morphic continued fractions}, Journal of Number Theory \textbf{91} (2001), no.\ 1, 39--66.

\bibitem{Bernstein1994} Daniel J.\ Bernstein, \emph{A noniterative 2-adic statement of the $3N+1$ conjecture}, Proceedings of the American Mathematical Society \textbf{121} (1994), no.\ 2, 405--408.

\bibitem{BernsteinLagarias1996} Daniel J.\ Bernstein and Jeffrey C.\ Lagarias, \emph{The $3x+1$ conjugacy map}, Canadian Journal of Mathematics \textbf{48} (1996), no.\ 6, 1154--1169.

\bibitem{BHZ2006} Valérie Berthé, Charles Holton, and Luca Q.\ Zamboni, \emph{Initial powers of Sturmian sequences}, Acta Arithmetica \textbf{122} (2006), no.\ 4, 315--347.

\bibitem{Dubickas2009} Artūras Dubickas, \emph{On integer sequences generated by linear maps}, Glasgow Mathematical Journal \textbf{51} (2009), no.\ 2, 243--252.

\bibitem{Lagarias1985} Jeffrey C.\ Lagarias, \emph{The $3x+1$ problem and its generalizations}, The American Mathematical Monthly \textbf{92} (1985), no.\ 1, 3--23.

\bibitem{LopezStoll2009} Josefina López and Peter Stoll, \emph{The $3x+1$ conjugacy map over a Sturmian word}, Integers \textbf{9} (2009), no.\ A13, 141--162.

\bibitem{LopezStoll2021} Josefina López and Peter Stoll, \emph{The $3x+1$ periodicity conjecture in $\mathbb{R}$}, 2021.

\bibitem{Lothaire2002} M.\ Lothaire, \emph{Algebraic Combinatorics on Words}, Encyclopedia of Mathematics and its Applications, vol.\ 90, Cambridge University Press, 2002, Chapter 2 (Sturmian words, by J.\ Berstel and P.\ Séébold). Cited here as the standard reference for the combinatorics-on-words background of complementary symmetric Rote sequences; the primary source defining this class (G.\ Rote, 1994) should be verified directly against primary literature before formal submission, as this repository's record does not independently confirm its exact publication details.

\bibitem{Phase2Abstract} \repo{theorems/phase2/ABSTRACT\_IRRATIONALITY\_CRITERION.md}, periodicity-conjecture-bridge repository, commit \texttt{a4ac46a}.
\bibitem{Phase2Surplus} \repo{theorems/phase2/SURPLUS\_THEOREM.md} and \repo{theorems/phase1/SURPLUS\_INVARIANT.md}, ibid.
\bibitem{Phase1Discrepancy} \repo{theorems/phase1/DISCREPANCY\_HEIGHT\_THEOREM.md}, ibid.
\bibitem{Phase1Repetition} \repo{theorems/phase1/REPETITION\_HEIGHT\_CRITERION.md}, ibid.
\bibitem{Phase2Discrepancy} \repo{theorems/phase2/ROTE\_ROOT\_DISCREPANCY.md}, ibid.
\bibitem{Phase2Transfer} \repo{theorems/phase2/ROTE\_TRANSFER\_THEOREM.md}, ibid.
\bibitem{Phase5Transfer} \repo{theorems/phase5/TRANSFER\_REDERIVATION.md}, ibid.
\bibitem{Phase3Automaton} \repo{theorems/phase3/WEIGHT\_PARITY\_CLASSIFICATION.md}, ibid.
\bibitem{Phase4Silver} \repo{theorems/phase4/SILVER\_EXTREMALITY.md}, ibid.
\bibitem{Phase4SilverAudit} \repo{theorems/phase4/SILVER\_CONSTANT\_AUDIT.md}, ibid.
\bibitem{Phase4Required} \repo{theorems/phase4/REQUIRED\_EXPONENT.md}, ibid.
\bibitem{Phase5Hostile} \repo{theorems/phase5/HOSTILE\_REFEREE\_REPORT.md}, ibid.
\bibitem{Phase6Arrow} \repo{theorems/phase6/ARROW\_BY\_ARROW\_RECONSTRUCTION.md}, ibid.
\bibitem{Phase4Quadratic} \repo{theorems/phase4/QUADRATIC\_CS\_ROTE\_THEOREM.md}, ibid.
\bibitem{Phase6Bounded} \repo{theorems/phase6/BOUNDED\_TYPE\_THEOREM.md}, ibid.
\bibitem{Phase4Chain} \repo{theorems/phase4/FINAL\_PROOF\_CHAIN.md}, ibid.
\bibitem{BHZGeneralMap} \repo{theorems/phase7/BHZ\_GENERAL\_INTERCEPT\_MAP.md}, ibid.
\bibitem{Phase7Verdict} \repo{theorems/phase7/PHASE7\_VERDICT.md}, ibid.
\bibitem{Phase7Sufficient} \repo{theorems/phase7/INTERCEPT\_SUFFICIENT\_CONDITION.md}, ibid.
\bibitem{Phase7Surplus} \repo{theorems/phase7/SURPLUS\_DECOMPOSITION.md} and \repo{theorems/phase7/INTERCEPT\allowbreak\_EXCESS.md}, ibid.
\bibitem{Phase7Flexible} \repo{theorems/phase7/FLEXIBLE\_MARGIN\_BOUNDARY.md}, ibid.
\bibitem{Phase7Structured} \repo{theorems/phase7/STRUCTURED\_INTERCEPTS.md}, ibid.
\bibitem{Phase7Drift} \repo{theorems/phase7/DRIFT\_COORDINATES.md}, ibid.
\bibitem{Phase7DriftDigits} \repo{theorems/phase7/DRIFT\_VS\_DIGITS.md}, ibid.
\bibitem{Phase7Intrinsic} \repo{theorems/phase7/INTRINSIC\_SQUARE\_ROUTE.md}, ibid.
\bibitem{Phase2Rotation} \repo{theorems/phase2/ROTATION\_EXAMPLE\_THEOREM.md}, ibid.
\bibitem{Phase1Gap} \repo{theorems/phase1/FULL\_PERIODICITY\_GAP.md}, ibid.

\end{thebibliography}

\end{document}
